% TOP_PROBLEM: {"id": 30006966, "title": "Liouville problem for stationary Navier–Stokes D-solutions in R^3", "category_id": 9, "category": {"id": 9, "name": "pde", "display_name": "Partial Differential Equations", "description": "PDEs and their applications in physics and geometry.", "slug": "pde", "order_index": 9, "created_at": "2026-07-31T15:26:25.671Z"}, "status": "open"}
% FIRST_POSTED: 2026-09-27
% !TeX program = lualatex
% ATTEMPT_STATUS: unresolved
% Model: GPT-6 Astra Ultra; continuation dated 27 September 2026.
\documentclass[11pt]{article}
\usepackage[margin=25mm]{geometry}
\usepackage{fontspec}
\setmainfont{Latin Modern Roman}
\usepackage{amsmath,amssymb,mathtools}
\usepackage{unicode-math}
\setmathfont{Latin Modern Math}
\usepackage{xurl,hyperref}
\hypersetup{colorlinks=true,urlcolor=blue}
\setcounter{secnumdepth}{0}
\setlength{\parindent}{0pt}
\setlength{\parskip}{5pt}
\setlength{\emergencystretch}{3em}
\begin{document}
\section*{\texorpdfstring{Liouville problem for stationary Navier–Stokes D-solutions in R\textasciicircum{}3}{Liouville problem for stationary Navier–Stokes D-solutions in R\textasciicircum{}3}}\label{30006966}
\subsection{Definitions and mathematical statement}
Let $u:{\mathbb{R}}^3\to{\mathbb{R}}^3$ and $p:{\mathbb{R}}^3\to{\mathbb{R}}$ be smooth and satisfy the stationary unforced system with viscosity normalized to one:
\[
 -\Delta u+(u\cdot\nabla)u+\nabla p=0,
 \qquad\nabla\cdot u=0.
\]
Assume finite Dirichlet integral and uniform vanishing at infinity:
\[
 \int_{{\mathbb{R}}^3}|\nabla u|^2{\,\mathrm{d}} x<\infty,
 \qquad \lim_{R\to\infty}\sup_{|x|\ge R}|u(x)|=0.
\]
Must $u\equiv0$? The finite-energy condition here is on the gradient, not the kinetic energy $\int|u|^2$, which is not an extra assumption. If the conclusion holds, the equation also makes the pressure constant. The spatial domain has no obstacle or physical boundary supplying energy.
\par\smallskip\textit{Formulation and concept references:} \href{https://arxiv.org/html/2506.14533v1}{[S1]}.

\subsection{Short English statement}
Must every smooth decaying stationary Navier--Stokes flow on all of three-dimensional space vanish when its gradient has finite square integral?
\subsection{Research attempt}

\paragraph{Scope, source audit, and attack (27 September 2026).}
This GPT-6 Astra Ultra attempt concerns stationary solutions on all of
$\mathbb R^3$, with the two hypotheses printed above. In particular,
$\nabla u\in L^2$ is given, whereas $u\in L^2$ is not. There is no
boundary on which one can prescribe a zero energy flux. The conclusion
must hold without a quantitative rate in the uniform decay assumption.

The supplied primary source [S1], Introduction and Theorems 1.1--1.2,
describes the $D$-solution problem and proves results with additional
control on a velocity antiderivative or on line integrals. Those are
conditional hypotheses, not consequences of the definition of a
$D$-solution. Its introduction also identifies the classical sufficient
condition $u\in L^{9/2}$, which is rederived below as a benchmark.
Chae's 2026 paper likewise proves conditional results involving head
pressure decay and records the remaining three-dimensional difficulty
[E1]. No unconditional result is inferred from either title.

The route investigated here is the direct energy argument: normalize the
pressure, retain every term in the cutoff identity, and determine exactly
which terms are controlled by the stated assumptions. The proof track
reaches a complete theorem with the extra $L^{9/2}$ hypothesis. The
counterexample track constructs a smooth divergence-free field satisfying
the functional hypotheses but violating that extra integrability. It is
explicitly checked not to solve the equation, so it invalidates a proposed
functional shortcut rather than the conjecture.

\paragraph{Step 1: what the Dirichlet hypothesis actually gives.}
\textbf{Lemma 487.1.} Under the stated hypotheses,
\[
 u\in L^6(\mathbb R^3)\cap L^\infty(\mathbb R^3),
 \qquad \|u\|_6\leq C\|\nabla u\|_2.
\]

\emph{Proof.}
Smoothness on compact sets and uniform decay outside large balls give
boundedness. For $\varepsilon>0$, the Sobolev function
$v_\varepsilon=(|u|-\varepsilon)_+$ has compact support, by uniform
decay. The chain rule gives
$|\nabla v_\varepsilon|\leq|\nabla u|$ almost everywhere. Apply the
ordinary three-dimensional Sobolev inequality to its smooth
approximations to obtain
$\|v_\varepsilon\|_6\leq C\|\nabla u\|_2$.
Let $\varepsilon\downarrow0$ and use monotone convergence.
This also explains why a nonzero additive constant, permitted by a
homogeneous gradient norm alone, is excluded here.\hfill$\square$

It follows by boundedness that $u\in L^q$ for every $q\geq6$, since
$\int|u|^q\leq\|u\|_\infty^{q-6}\int|u|^6$.
There is no reverse inclusion into $L^{9/2}$ on a space of infinite
measure. The tempting interpolation with $L^2$ would require an
assumption absent from the problem.

We use the standard $L^q$ boundedness of the Riesz transforms:
$\|R_iR_j f\|_q\leq C_q\|f\|_q$ for $1<q<\infty$.
The multiplier is $-\xi_i\xi_j/|\xi|^2$. This is the
Calder\'on--Zygmund estimate used for pressure in [E1], equation (2.30).
Its boundedness is a standard analytic input, not a new assertion proved
in this notebook. All exponents at which it is used are strictly between
one and infinity.

\paragraph{Step 2: pressure normalization without a hidden decay assumption.}
Define
\begin{equation}\label{ra487:pressure}
 p_0=\sum_{i,j=1}^3R_iR_j(u_i u_j).
\end{equation}
Since $u\in L^6$, this is a well-defined $L^3$ function and
\[
 \|p_0\|_3\leq C\|u\|_6^2.
\]
One must still check that the pressure appearing in the equation differs
from $p_0$ by only a constant; simply solving the pressure Poisson
equation leaves a possible harmonic term.

\textbf{Lemma 487.2.} The given smooth pressure satisfies $p=p_0+c$
as a distribution, for a constant $c$.

\emph{Proof.}
Taking divergence of the stationary equation and using
$\nabla\cdot u=0$ gives
\[
 -\Delta p=\sum_{i,j}\partial_i\partial_j(u_i u_j)
                         =-\Delta p_0.
\]
Let $w=\nabla(p-p_0)$. The equation also writes
\[
 w=\Delta u-\operatorname{div}(u\otimes u)-\nabla p_0.
\]
The right side is tempered because $u\in L^6$ and $u\otimes u,p_0\in
L^3$. Moreover $\Delta w=0$. A tempered harmonic distribution is a
harmonic polynomial: its Fourier transform is supported at the origin
and is a finite linear combination of derivatives of the delta
distribution. We use this basic distribution fact for each component.

For any $\phi\in C_c^\infty(\mathbb R^3)$, put
$\phi_R(x)=R^{-3}\phi(x/R)$. Distributional differentiation and H\"older
give
\begin{align*}
 |\langle\Delta u,\phi_R\rangle|
 &\leq C_\phi R^{-5/2}\|u\|_6,\\
 |\langle\operatorname{div}(u\otimes u),\phi_R\rangle|
   +|\langle\nabla p_0,\phi_R\rangle|
 &\leq C_\phi R^{-2}(\|u\|_6^2+\|p_0\|_3).
\end{align*}
For example, the $L^{6/5}$ norm of $\Delta\phi_R$ is
$R^{-5+3(5/6)}\|\Delta\phi\|_{6/5}=R^{-5/2}\|\Delta\phi\|_{6/5}$;
the $L^{3/2}$ norm of $\nabla\phi_R$ is $R^{-2}$ times a fixed
constant. Thus $\langle w,\phi_R\rangle\to0$ for every such $\phi$.

If a component of the polynomial $w$ were nonzero, pairing it with
$\phi_R$ would be a polynomial in $R$. The coefficients are its
homogeneous polynomial pieces paired with $\phi$. The limit zero for
every $\phi$ forces every coefficient to vanish, including the constant
one. Hence $w=0$. A distribution with zero gradient on connected
$\mathbb R^3$ is constant, proving the claim.\hfill$\square$

This argument does not assume that a general smooth $L^3$ function has a
uniform pointwise limit at infinity. That implication would be false
without further control of possible narrow peaks. The normalization
needed below is the integral normalization $p_0\in L^3$.

\paragraph{Step 3: the exact cutoff identity and its unresolved term.}
Fix $\phi\in C_c^\infty(\mathbb R^3)$ with $0\leq\phi\leq1$,
$\phi=1$ on the unit ball and support in the ball of radius two.
Write $\phi_R(x)=\phi(x/R)$ in this paragraph, and
$A_R=\{R<|x|<2R\}$. This cutoff differs from the normalized test
function in Lemma 487.2. Multiplying the equation by $u\phi_R$ and
integrating gives
\begin{equation}\label{ra487:energy}
 \int |\nabla u|^2\phi_R
 =\frac12\int |u|^2\Delta\phi_R
     +\int\left(\frac{|u|^2}{2}+p_0\right)u\cdot\nabla\phi_R.
\end{equation}
All terms have compact support or a compactly supported derivative, so
these integrations require no unproved global integrability of $|u|^2$.
For the viscous term use
$\nabla(|u|^2/2)=u\cdot\nabla u$. For convection use
$(u\cdot\nabla)u\cdot u=u\cdot\nabla(|u|^2/2)$ and integrate the
divergence. The pressure term integrates to
$-\int p\,u\cdot\nabla\phi_R$; its constant part vanishes because
$\int u\cdot\nabla\phi_R=0$.

Set $U_R=\|u\|_{L^6(A_R)}$ and $P_R=\|p_0\|_{L^3(A_R)}$.
The first term is controlled by
\begin{equation}\label{ra487:diffusion}
 R^{-2}\int_{A_R}|u|^2
 \leq C R^{-2}|A_R|^{2/3}U_R^2
 \leq C U_R^2\longrightarrow0.
\end{equation}
The two remaining absolute-value estimates are
\begin{align}
 R^{-1}\int_{A_R}|u|^3
       &\leq C R^{1/2}U_R^3,\label{ra487:convection}\\
 R^{-1}\int_{A_R}|p_0||u|
       &\leq C R^{1/2}P_R U_R.\label{ra487:pressureflux}
\end{align}
Here the residual volume exponent is
$1-3/6=1/2$ for the first integral and
$1-1/3-1/6=1/2$ for the second. Both $U_R$ and $P_R$ tend to zero,
but no rate has been proved that defeats the factor $R^{1/2}$.

This is the exact point where testing by $u$ without a cutoff would
conceal a gap. Formally discarding the last integral in
\eqref{ra487:energy} amounts to assuming that the energy flux through
large scales vanishes. The given gradient energy makes the left side
converge to the total Dirichlet integral by dominated convergence; it
does not, by itself, justify discarding that flux.

There is a precise conditional criterion even at the present exponents:
if a sequence $R_j\to\infty$ satisfies
\begin{equation}\label{ra487:tailcriterion}
 R_j^{1/2}\bigl(U_{R_j}^3+P_{R_j}U_{R_j}\bigr)\longrightarrow0,
\end{equation}
then \eqref{ra487:energy} implies $\nabla u=0$, and uniform decay
implies $u=0$. This is a proved sufficient condition, not an established
property of every $D$-solution. Absolute estimates are only one possible
way to prove zero flux; a successful argument could instead exploit
cancellation in the signed integral.

\paragraph{A complete benchmark: the extra $L^{9/2}$ hypothesis.}
\textbf{Proposition 487.3.} If the solution in the problem also satisfies
$u\in L^{9/2}(\mathbb R^3)$, then $u\equiv0$ and $p$ is constant.

\emph{Proof.}
The Riesz transform estimate now gives
$p_0\in L^{9/4}$ with
$\|p_0\|_{9/4}\leq C\|u\|_{9/2}^2$.
The definitions of the Riesz transform on $L^3\cap L^{9/4}$ agree,
by approximation with smooth compactly supported functions, so this is
the same pressure as in Lemma 487.2.
Writing $V_R=\|u\|_{L^{9/2}(A_R)}$ and
$Q_R=\|p_0\|_{L^{9/4}(A_R)}$, H\"older gives
\begin{align*}
 R^{-1}\int_{A_R}|u|^3
  &\leq R^{-1}|A_R|^{1/3}V_R^3\leq C V_R^3,\\
 R^{-1}\int_{A_R}|p_0||u|
  &\leq R^{-1}|A_R|^{1-4/9-2/9}Q_R V_R
       \leq C Q_R V_R.
\end{align*}
Both terms tend to zero because the global integrals defining these
norms are finite. The diffusion term tends to zero by
\eqref{ra487:diffusion}. Let $R\to\infty$ in
\eqref{ra487:energy}; it follows that $\int|\nabla u|^2=0$.
Smoothness and connectedness imply that $u$ is constant, and decay makes
the constant zero. Substitution in the equation gives $\nabla p=0$.
\hfill$\square$

The special exponent comes from the elementary expression
\[
 R^{-1}|A_R|^{1-3/q}\asymp R^{2-9/q}.
\]
It has exponent zero at $q=9/2$, positive exponent at $q=6$, and
negative exponent for $3\leq q<9/2$. With the corresponding pressure
space $L^{q/2}$, the same proof works in that latter range. This is
criticality for this particular cutoff estimate; it should not be confused
with the scaling-critical Navier--Stokes velocity space $L^3$.
The $L^{9/2}$ theorem is known, as explicitly recorded in [S1]; the
purpose of the reconstruction is to locate exactly the missing decay.

\paragraph{A functional counterexample to the proposed integrability upgrade.}
For $1/2<\alpha\leq2/3$, put $\beta=(\alpha+1)/2$ and define
\begin{equation}\label{ra487:swirl}
 U_\alpha(x)=(1+|x|^2)^{-\beta}(-x_2,x_1,0).
\end{equation}
This field is smooth everywhere. Its divergence vanishes: the radial
gradient of $(1+|x|^2)^{-\beta}$ is orthogonal to $(-x_2,x_1,0)$,
whose own divergence is zero. Also $|U_\alpha(x)|\leq C(1+|x|)^{-\alpha}$,
so it vanishes uniformly. Direct differentiation gives
$|\nabla U_\alpha(x)|\leq C(1+|x|)^{-\alpha-1}$ and consequently
\[
 \int_{|x|>1}|\nabla U_\alpha|^2
 \leq C\int_1^\infty r^{-2\alpha}\,\mathrm dr<\infty.
\]
On a fixed angular belt away from the $x_3$ axis,
$|U_\alpha|$ is bounded above and below by positive constants times
$r^{-\alpha}$. In fact the angular factor is $\sin\theta$ and its
$q$th power has a positive finite integral over the sphere. Hence
\begin{equation}\label{ra487:swirlintegrability}
 U_\alpha\in L^q(\mathbb R^3)
        \quad\Longleftrightarrow\quad q\alpha>3
                 \qquad(q>0).
\end{equation}
For $\alpha=3/5$, this gives $U_\alpha\in L^6$ but
$U_\alpha\notin L^{9/2}$ and $U_\alpha\notin L^2$.
Moreover
\[
 R^{-1}\int_{A_R}|U_{3/5}|^3\asymp R^{1/5}.
\]
Thus even smoothness, divergence freedom, uniform decay, finite Dirichlet
integral and the resulting $L^6$ bound together cannot force the
absolute convection estimate to vanish.

This is not a stationary Navier--Stokes counterexample. To verify that
point rather than leave it implicit, write $f(r)=(1+r^2)^{-\beta}$.
Then
\[
 \Delta U_\alpha=g(r)(-x_2,x_1,0),\qquad
 g(r)=f''(r)+\frac4r f'(r)
 =-2\beta(1+r^2)^{-\beta-2}
       \bigl(5+(2-\alpha)r^2\bigr),
\]
where the expression at zero is understood by continuity. Also
$(U_\alpha\cdot\nabla)U_\alpha=-f(r)^2(x_1,x_2,0)$.
If some smooth pressure made the stationary equation hold, its gradient
would equal
$g(r)(-x_2,x_1,0)+f(r)^2(x_1,x_2,0)$.
The integral of this vector field around a circle of radius $r>0$ in
the plane $x_3=0$ is $2\pi r^2g(r)$, which is nonzero. A gradient has
zero integral around a closed curve, a contradiction.

For radial cutoffs the signed convection term of this example is
actually zero because $U_\alpha\cdot\nabla\phi_R=0$. Thus the example
has a second lesson: absolute estimates may lose essential cancellation.
It disproves a functional implication used by a naive proof; it does not
show that the true stationary energy flux is nonzero for an actual
$D$-solution.

\paragraph{Attempt through head pressure: what the sign does and does not give.}
Let $Q=p_0+|u|^2/2$ and $\omega=\nabla\times u$. A calculation using
the equation and the pressure Poisson identity yields
\begin{equation}\label{ra487:bernoulli}
 -\Delta Q+u\cdot\nabla Q=-|\omega|^2.
\end{equation}
Indeed $\Delta(|u|^2/2)=|\nabla u|^2+u\cdot\nabla Q$, while
$\Delta p_0=-\sum_{i,j}\partial_i u_j\partial_j u_i$; subtracting
gives $\Delta Q-u\cdot\nabla Q=|\omega|^2$.
This identity motivates head-pressure methods such as [E1], but does
not by itself eliminate the right side.

Even if one has additionally justified the pointwise normalization
$Q(x)\to0$, the maximum principle gives only $Q\leq0$. For example,
the nonzero scalar function $q(x)=-(1+|x|^2)^{-1/2}$ satisfies
\[
 q(x)\to0,\qquad -\Delta q=-3(1+|x|^2)^{-5/2}<0.
\]
It is therefore false that a decaying nonpositive solution of an equation
of the form $-\Delta q=-f$, $f\geq0$, must vanish. This scalar test is
not asserted to arise from a Navier--Stokes pair; it pinpoints the missing
information in using the maximum principle alone. A proof must couple
\eqref{ra487:bernoulli} to the velocity equation and obtain an additional
integral or level-set estimate. No such unconditional estimate is
established here.

\paragraph{Verification and outcome.}
The companion exact-arithmetic script records the exponents
$2-9/q$, the integrability thresholds $q\alpha>3$, and the coefficients
in $g(r)$ for rational test parameters, including $\alpha=3/5$.
Those finite algebra checks help catch a wrong sign or scale; they are
not numerical evidence for an infinite-domain Liouville theorem.
The mathematical verification above separately checks pressure constants,
compact support of each test, the zero mode in the harmonic pressure
argument, and the difference between a manufactured field and a solution.

\textbf{Outcome: UNRESOLVED.} The proved results are the $L^6$ consequence
of the hypotheses, the canonical pressure normalization, the exact energy
identity, the tail criterion \eqref{ra487:tailcriterion}, and the known
$L^{9/2}$ sufficient condition. The functional counterexample explains
why an integrability upgrade cannot be made without using further content
of the stationary equation.

The first unresolved step for this route is proving, from the original
hypotheses alone, that the signed flux
$\int (|u|^2/2+p_0)u\cdot\nabla\phi_R$ tends to zero along some
sequence of radii. A stronger sufficient goal is
\eqref{ra487:tailcriterion}; another is $u\in L^{9/2}$. Neither is
proved for arbitrary $D$-solutions. Concrete continuations are to exploit
cancellation between pressure and convection, to control low-frequency
velocity tails using the equation, or to obtain a head-pressure level-set
estimate strong enough to close \eqref{ra487:bernoulli}. None has been
completed here, and no counterexample satisfying the stationary system
has been produced.

\subsection{Sources}
\begin{itemize}
\item[S1]Coiculescu and Jincheng Yang, Conditional Liouville Theorems for the Navier–Stokes Equations, arXiv v1 (2025).
\url{https://arxiv.org/html/2506.14533v1}.
\textit{Formulation source.}
\item[E1] D. Chae, \emph{Liouville Type Theorems for the Stationary
Navier--Stokes Equations in $\mathbb R^3$}, Communications in Mathematical
Physics 407 (2026), article 53, Introduction and equation (2.30).
\url{https://link.springer.com/article/10.1007/s00220-026-05555-y}.
\textit{Primary published source, 5 February 2026; conditional results and
pressure estimates verified 27 September 2026.}
\end{itemize}
\end{document}
